% Limitations and improvements
\section{Limitations and improvements}\label{sec:limitations}

\paragraph{Final accuracy against the target.}
The brief sets a target of up to \qty{99}{\percent} mean accuracy. We reach \GateMeanAcc{} at gate
level, \FreshMeanAcc{} on fresh corners, and \CircuitMeanAcc{} on the test circuit, so the target
is not met, by about one to two percentage points. The analysis of Section~\ref{sec:error} locates
the gap precisely: the composition is exact (SPICE oracle \OracleMeanAcc, circuit composition from
SPICE gate values \CompSpiceMeanAcc), so all of the remaining error comes from the transistor models,
and the single largest contribution, \ErrorSourceFirst{}, would on its own raise the gate accuracy
to \ErrorSourceFirstAcc{} if it were exact. The limitations below are grouped by part, each with
the change we would make.

\begin{itemize}
  \item \textbf{Data and sampling}
  \begin{itemize}
    \item \emph{ON transistors at intermediate $|V_{DS}|$.} The sampling targets ON transistors at
          $V_{DS}\approx 0$ and OFF transistors at the rails, but an ON NMOS whose source floats at
          an internal node (the upper NMOS of NAND2~\texttt{10}, $x \approx
          \VxOneZeroFrac\,V_{DD}$) sits at a substantial $|V_{DS}|$ near threshold and causes the
          worst gate case. \emph{Improvement:} an ON-case analogue of source pinning, with the
          source drawn near $V_{DD} - V_{th}$ (NMOS) or $V_{th}$ (PMOS) and the drain on the rail;
          like the existing rules it follows from circuit physics, not from evaluation data.
    \item \emph{Dataset size.} The learning curve (Figure~\ref{fig:lc}) still slopes down at
          \LCMaxSize{} training rows. \emph{Improvement:} more rows, concentrated in the
          oversampled regions where the error is largest.
    \item \emph{Junction current floor.} The OFF currents sit on a constant floor of about
          $10^{-14}$\,A at low temperature (Figure~\ref{fig:trends}), consistent with BSIM4's
          default junction current when the netlist gives no junction areas, so junction leakage
          does not scale with $W$. \emph{Improvement:} give every transistor \texttt{AS}/\texttt{AD}
          and \texttt{PS}/\texttt{PD} from its width in all decks.
  \end{itemize}
  \item \textbf{Transistor models}
  \begin{itemize}
    \item \emph{PMOS gate tunnelling.} The largest error source at gate level is the ON PMOS gate
          current (Table~\ref{tab:errorsources}), whose error depends on the oxide thickness and
          does not average out over a corner. \emph{Improvement:} a dedicated model for $I_G$ with
          the physical $\exp(-\alpha\, t_{ox})$ dependence built into the features or the target.
    \item \emph{Accuracy floor at tiny currents.} Currents that cross zero (ON $I_D$, $I_S$) or are
          close to the \IFloor{} floor (OFF $I_S$) have unbounded relative error, and below about
          $10^{-14}$\,A even the SPICE reference is only good to a per cent or two
          (\SpiceFloorRel{} of the smallest gate leakage). \emph{Improvement:} train on the
          quantities that enter $I_{VDD}$ (the supply-side terminal currents) with a loss weighted
          by their share of the gate leakage, instead of on each terminal current separately.
    \item \emph{Grid edge and model size.} All four chosen GBDTs use the largest number of leaves
          in the grid (\ChosenHyperNmosOff), so larger trees may still help, and the linear-leaf
          models are large files. \emph{Improvement:} extend the grid, then prune or distil the
          final models.
  \end{itemize}
  \item \textbf{Gates}
  \begin{itemize}
    \item \emph{SPICE node voltages are still required.} As the brief prescribes, the gate
          prediction reads the internal and output node voltages from SPICE, so it is not yet a
          SPICE-free leakage estimator; the optional ML node solve was not implemented.
          \emph{Improvement:} solve KCL at the internal nodes with the ML currents (a 1-D root
          search per node, coupled for AND2/OR2) and compare the solved voltages with SPICE.
    \item \emph{Heuristic KCL threshold.} ON transistors below \KclVdsMax{} get their channel
          current from KCL. This works because the gap is wide (in-gate values at most \OnVdsMax,
          the smallest genuinely biased ON transistor at \OnBiasedVdsMin), but a gate with an
          intermediate $|V_{DS}|$ would fall between the two regimes. \emph{Improvement:} solve the
          channel currents of all ON transistors consistently with the node voltages, which the
          ML node solve above would provide.
    \item \emph{Ideal gate voltages.} The models only know $V_G\in\{0, V_{DD}\}$; the AND2/OR2
          inverter input is snapped (at most \GateSnapMax{} here). In circuits with degraded logic
          levels this would no longer be negligible. \emph{Improvement:} sample $V_G$ in a narrow
          band around each rail.
  \end{itemize}
  \item \textbf{Test circuit}
  \begin{itemize}
    \item \emph{Loading term.} The loading term assumes ideal $0$/$V_{DD}$ gate inputs and splits a
          gate's $I_{in}$ evenly over its high pins. For the NAND2-only parity checker this is
          essentially exact (\CompSpiceMeanAcc{} with SPICE gate values), but for gates with
          asymmetric inputs, such as the AND2/OR2 first stage, the even split is an approximation.
          \emph{Improvement:} predict $I_{in}$ per input pin in the gate model.
    \item \emph{Systematic error accumulates.} Corner-dependent model errors are common to all
          gates of a corner and add up in the circuit (Section~\ref{sec:errcircuit}), so the
          circuit accuracy is bounded by the systematic part of the gate error. \emph{Improvement:}
          the transistor-model improvements above; a larger circuit would not average this error
          away.
  \end{itemize}
\end{itemize}
